\chapter{Connections on vector bundles} \label{ch_connections}
In this chapter, we develop the theory of connections over complex vector bundles. The notion of a connection itself is familiar from riemannian geometry, and the definition can be adapted to abstract vector bundles with minor modifications. The key idea behind connections is to provide a way of differenciating sections of vector bundles in a coordinate-independent way. The word connection is used with different meanings in the literature (all of them related of course), so we devote a part of the chapter to study these definitions. We also define the curvature of a connection from different points of view. The first definition we give should be familiar from riemannian geometry, but we will introduce a new way to think about curvatures as \(\mathrm{End}(E)\)-valued forms. Once these basic concepts are introduced, we give a proof of the existence of the Chern connection; a connection that is simultaneously compatible with a hermitian inner product and the complex structure. We also define the First Real Chern class and study some properties of the curvature of the Chern connection.

\Definition{}{Let \(\pi \colon E \to M\) be a smooth complex vector bundle of rank \(r\) over a smooth manifold \(M\). A \textit{connection}\index{Connection!on a vector bundle} on \(E\) is a map
\[
\begin{array}{cccc}
    \nabla \colon & \Gamma(T_\C M) \times \Gamma(E) & \longrightarrow & \Gamma(E) \\
    &\ \ \ \ \ \ \ \ (X, \sigma) & \mapsto & \nabla_X \sigma  
\end{array}
\]
satisfying the following properties:
\begin{enumerate}
    \item \(\nabla\) is \(\mathscr{C}^\infty(M, \C)\)-linear on its first component.
    \item For every \(f \in \mathscr{C}^\infty(M;\C)\), and \(\sigma \in \Gamma(E)\),
    \[
    \nabla_X (f \cdot \sigma) = X(f) \cdot \sigma + f \cdot \nabla_X \sigma.
    \]
\end{enumerate} \label{def_connection}
}
The expression \(\nabla_X \sigma\) is usually called the \textit{covariant derivative of \(\sigma\) in the direction of \(X\)}\index{Covariant derivative}. One can show that \((\nabla_X \sigma)_p\) for \(p \in M\) depends only on the value of \(X\) at \(p\) and the value of \(\sigma\) in an arbitrarily small neighborhood of \(p\), so for any open subset \(U \subset M\), restriction of the connection defines a connection on the restriction of the bundle \(E|_U\) (which is a smooth manifold of its own). See \parencite[Prop. 4.3]{lee_introRM} for a proof for connections on the real tangent bundle of a manifold, which applies almost verbatim to this case. \\

Sometimes, one is given the following definition of a connection:

\Definition{}{Let \(\pi \colon E \to M\), a \textit{connection} on \(E\) is a \(\C\)-linear map \(\nabla\colon \Gamma(E) \to \Gamma(T^*_\C M \otimes E)\) such that for every \(f \in \mathscr{C}^\infty(M;\C)\), \(\sigma \in \Gamma(E)\), we have 
\begin{equation} \label{def_connection_eq}
\nabla (f \cdot \sigma)=\diff f \otimes \sigma+f \cdot \nabla \sigma
\end{equation}
\label{def_tot_cov_der}}

The last definition is sometimes called a \textit{total covariant derivative}\index{Total covariant derivative}. Observe that a connection in the sense of definition \ref{def_connection} induces a total covariant derivative: for any \(\sigma \in \Gamma(E)\), we define \(\nabla \sigma \in \Gamma(T^*M \otimes E)\) to be the section such that \(\nabla\sigma(X)=\nabla_X \sigma\) (we are using here the vector bundle isomorphism \(T^*_\C M \otimes E \cong \mathrm{Hom}(T_\C M , E)\)). Reciprocally, a total covariant derivative induces a connection by simply defining \((X, \sigma) \mapsto \nabla\sigma (X)=:\nabla_X \sigma\). By an abuse of notation, we will refer to both constructions by the name connection. When we consider real vector bundles, we may still define connections analogously (we do not consider the complexified tangent and cotangent bundles in that case). Moreover, observe that \(\Omega_\C^1(M;E) \cong \Gamma(\bigwedge ^1T^*_\C M \otimes E)=\Gamma(T^*_\C M \otimes E)\), hence a connection is a map \(\nabla \colon \Omega_\C^0(M;E) \to \Omega_\C^1(M;E)\) satisfying the ``Leibniz rule'' of definition \ref{def_tot_cov_der}. This is the definition one can find in \parencite{huybrechts2005complex} and \parencite{kobayashi_complex_vector_bundles}.

\begin{obs}[\textbf{Coordinate expression of connections}] \label{coordinate_exp_connection}
    Consider a local frame \((s_1, \ldots s_r)\) of a complex vector bundle \(E\) defined on some open set \(U \subset M\). For each \(1 \leq j \leq r\), \(\nabla s_j \in \Omega_\C^1(E)\), and hence it may be written as
    \begin{equation} \label{def_connection_1_form}
    \nabla s_j = \omega^i_j \otimes s_i
    \end{equation}
    for some \(\omega^i_j \in \Omega_\C^1(M)\) that are called \textit{connection one-forms}. If we consider a section \(\xi=\xi^j\cdot  s_j \in \Gamma(E)\), then the coordinate expression of its image is
    \[
    \nabla \xi = \nabla(\xi^j s_j) = \diff \xi^j \otimes s_j + \xi^j \cdot  \nabla s_j = \parent{\diff \xi^k + \xi^j \cdot \omega^k_j} \otimes s_k.
    \]
    Equivalently,
    \[
    \nabla_X \xi = (X(\xi^k)+\xi^k \cdot \omega^k_j(X)) \cdot s_k.
    \]
    A section \(\xi\) is said to be \textit{parallel}\index{Parallel section} if \(\nabla \xi \equiv 0\). Observe that, with respect to the trivialization defined by the local frame \((s_1, \ldots s_r)\) (proposition \ref{prop_local_frames_associated_trivializations}), the section \(\xi\) can be identified with the smooth map \(M \to \C^r, m \mapsto (\xi^1(m), \ldots \xi^r(m))\). Then, \(\Phi(\nabla \xi)\) (extending \(\Phi\) to the tensor product \(T_\C^*M \otimes E\)) is given by \((\diff \xi^1+\xi^k\cdot \omega_k^1, \ldots \diff \xi^r + \xi^k \cdot \omega_k^r)\), which we may write as \(\diff \xi + \omega \cdot \xi\) for a matrix of one-forms \(\omega=(\omega^j_k)\). This is sometimes stated by saying that, locally, every connection is \(\diff+ A\) for some matrix of one-forms. \newline \newline

    Moreover, if we choose another local frame \((\tilde{s}_1, \ldots \tilde{s}_r)\) whose domain of definition overlaps with that of \((s_1, \ldots s_r)\), proposition \ref{prop_change_of_coordintes_frames} claims that there are smooth functions \(\tau^\ell_m\) such that \(\tilde{s}_k=\tau^j_k \cdot s_j\). Then, a straightforward calculation (see \parencite[Page 195]{lee_introCM}) shows that the connection one-forms with respect to the new frame relate to the original ones by
    \begin{eqnarray}\label{eq_transf_connection_one_form}
        \tilde{\omega}=(\tau^{-1})^p_j\cdot \diff \tau^j_k + (\tau^{-1})^p_\ell \cdot \tau^j_k \cdot \omega^\ell_j.
    \end{eqnarray}
    % It is sometimes useful to consider the matrix of \(1\)-forms \(\omega=(\omega^i_j)_{1 \leq i,j \leq r}\), because then all equations \eqref{def_connection_1_form} may be expressed together as \(\nabla s = s \otimes \omega\), where \(s:=(s_1, \ldots s_r)\), \(\nabla s\) is understood to act component-wise, and \(s \otimes \omega\) is understood as the usual product of a row vector and a matrix, but replacing ordinary products by tensor products.
    %  Consider now a different local frame \((\tilde{s}_1, \ldots \tilde{s}_r)\). In the open set where they overlap, there is a \(\mathrm{GL}(r, \C)\)-valued smooth map \(\tau\), such that \(\tilde{s}_p=s_p \cdot\tau(p)\) (this is just a matrix expression for proposition \ref{prop_change_of_coordintes_frames}), then, the connection form with respect to the new trivialization is
    % \[
    % \omega'=\tau^{-1}\cdot \omega \cdot \tau + \tau^{-1} \cdot \diff \tau,
    % \] 
    % where the exterior differential in meant to act component-wise. Indeed,
    % \[
    % \omega' = \nabla s' = \nabla s \cdot \tau = \nabla s \otimes  \tau + s\otimes  \diff \tau = s \otimes  \omega \cdot \tau + s \otimes \diff \tau = s' \otimes (\tau^{-1} \cdot  \omega \cdot  \tau + \tau^{-1} \cdot \diff \tau).
    % \]
\end{obs}

Finally, observe that non-trivial connections always exist. A quick way of obtaining one is the following: choose an open cover of \(M\) by domains of local frames \(\corch{U_\alpha}_{\alpha \in A}\), and let \(\corch{\varphi_\alpha}_{\alpha \in A}\) be a partition of unity subordinate to the cover. On each \(U_\alpha\) define a connection \(\nabla^\alpha\) by setting all connection one-forms to zero. Then, 
\begin{equation} \label{def_nontrivial_connection}
\nabla:=\sum_\alpha \varphi_\alpha \cdot \nabla^\alpha
\end{equation}
is a non-zero connection on \(E\).

\section{Curvature of a connection}
Every connection has an associated curvature form. One may find this object introduced in different equivalent ways in the literature. Some texts (\parencite{kobayashi_complex_vector_bundles} for instance) define the curvature by \(\nabla \circ \nabla\). If we take this expression and definition \ref{def_connection} literally, it does not make sense. Indeed, \(\nabla\) is not an endomorphism; it can not be composed with itself. To give this definition properly, one has to extend \(\nabla\) to a family of operators \(\diff_\nabla \colon \Omega_\C^p(M;E)\to \Omega_\C^{p+1}(M,E)\). This is done in theorem \ref{Th_ext_cov_der} bellow. However, it is quite useful to see the curvature as an \(\mathrm{End}(E)\)-valued two-form, and that is not immediate from the definition \(\diff_\nabla\circ \diff_\nabla\). To see the curvature that way, let us work with reference frames. Let \((s_1, \ldots s_r)\) be a smooth local frame for \(E\). Define the \textit{curvature two forms}\index{Curvature two forms (vector bundles)} by 
\begin{equation} \label{def_curv2_forms}
\Omega^k_j:=\diff \omega^k_j+ \omega^k_\ell\wedge \omega^\ell_j.
\end{equation}
Then, we define, locally, the curvature as
\begin{equation} \label{F_nabla_curv}
    F_\nabla = \Omega^k_j \otimes s_k \otimes t^j.
\end{equation}
, where \((t^1, \ldots t^r)\) is the dual frame to \((s_1, \ldots s_r)\). Since \(E \otimes E^*\cong \mathrm{End}(E)\), we may see \(F_\nabla\) as a (locally defined) \(\mathrm{End}(E)\)-valued two-form. However, if we choose another frame, and we use the transformation equation given by \eqref{eq_transf_connection_one_form}, a straightforward (but calculations) shows that \(F_\nabla=\tilde{\Omega}^k_j\otimes \tilde{s}_k\otimes \tilde{t}^j\). Hence, equation \eqref{F_nabla_curv} defines \(F_\nabla\) globally. We can now see how this relates to the operators \(\diff_\nabla\). 
% Whenever we have a connection, we also have a curvature form. Again, one can find different definitions of curvature in the literature, so we will give the most common definitions one may find and look at the way they are related.

% \Definition{\textbf{Curvature of a connection}}{Let \(\pi \colon E \to M\) be a smooth complex vector bundle over a smooth manifold \(M\), and suppose that it is endowed with a connection \(\nabla\). The \textit{curvature}\index{Curvature!of a connection on a vector bundle} of \(\nabla\) is the map \(R^\nabla \colon \Gamma(T_\C M ) \times \Gamma(T_\C M) \times \Gamma(E)\to \Gamma(E)\) defined by
% \[
% R^\nabla(X,Y)(\sigma)=\nabla_X \nabla_Y \sigma - \nabla_Y \nabla _X \sigma -\nabla_{[X,Y]} \sigma.
% \] 
% We say that a connection is \textit{flat}\index{Connection!flat} if \(R^\nabla=0\).\label{def_connection}
% }
% As we will show in the following lemma, \(R^\nabla\) is \(\mathscr{C}^\infty(M;\C)\)-linear in all three arguments, and it is antisymmetric with respect to its vector field arguments. Thus, it defines a section of the bundle \(\bigwedge^2 T_\C^*M \otimes E \otimes E^* \cong \bigwedge^2 T_\C^*M \otimes \mathrm{End}(E)\). As we detail in remark \ref{curvature_as_EndE_two_form}, this means we can see the curvature as an \(\mathrm{End}(E)\)-valued two-form on \(M\). A simple computation shows that \(R^\nabla\) is linear with respect to all of its components.

% \Lemma{}{Let \(\nabla\) be a connection on a complex vector bundle \(\pi\colon E \to M\). Then, its curvature is \(\mathscr{C}^\infty(M;\C)\)-linear with respect to all of its components.}
% \begin{proof}
%     Let us see first that it is linear in its first vector field entry (for the second the procedure is analogous).
%     \begin{align*}
%         R^\nabla(fX,Y) \sigma & = \nabla_{fX} \nabla_Y \sigma -\nabla_Y\nabla_{fX} \sigma -\nabla_{[fX, Y]}\sigma = f \cdot  \nabla_X\nabla_Y \sigma - \nabla_Y \parent{f \cdot \nabla _X \sigma} - \nabla_{f\cdot [X,Y]-Y(f) \cdot X} \sigma \\
%         & = f \cdot  \nabla_X\nabla_Y \sigma  - Y(f) \cdot \nabla_X \sigma -f \cdot \nabla_Y \nabla_X \sigma - f \cdot \nabla_{[X,Y]} \sigma + Y(f) \cdot \nabla_X \sigma \\
%         & = f \cdot \parent{\nabla_X\nabla_Y \sigma -\nabla_Y\nabla_X \sigma - \nabla_{[X,Y]} \sigma} = f \cdot R^\nabla(X,Y) \sigma.
%     \end{align*}
%     Now for the \(\Gamma(E)\)-component, we have
%     \begin{align*}
%         R^\nabla(X,Y)(f \cdot \sigma)  = & \nabla_X \nabla_Y (f \cdot \sigma) - \nabla_Y \nabla_X(f \cdot \sigma) - \nabla_{[X,Y]}(f \cdot \sigma) \\
%          = & \nabla_X(Y(f) \cdot \sigma+f \cdot \nabla_Y \sigma) - \nabla_Y (X(f) \cdot \sigma+ f \cdot \nabla_X \sigma) - [X,Y](f) \cdot \sigma - f\cdot \nabla_{[X,Y]} \sigma \\
%          = &  X(Y(f)) \cdot \sigma+ Y(f) \cdot \nabla_X \sigma+X(f)\cdot \nabla_Y \sigma+ f\cdot \nabla_X \nabla_Y \sigma-Y(X(f)) \cdot \sigma-X(f) \nabla_Y \sigma \\
%          &  -Y(f) \cdot \nabla_X \sigma -f \cdot \nabla_Y \nabla_X \sigma-[X,Y](f)\cdot \sigma - f\cdot \nabla_{[X,Y]} \sigma \\
%           = &f \cdot \parent{\nabla_X \nabla_Y\sigma-\nabla_Y \nabla_X \sigma- \nabla_{[X,Y]} \sigma} - (X(Y(f))-Y(X(f))-[X,Y](f)) \cdot \sigma \\=& f \cdot R^\nabla(X,Y) \sigma. 
%     \end{align*}
% \end{proof}

% \begin{obs}  \label{curvature_as_EndE_two_form}
%     Let us see in a very explicit manner how to see the curvature as en element of \(\Omega_\C^2(\mathrm{End}(E))\). We will do this with some level of detail because the kind of identification the curvature is an example of is quite common in the literature when working with vector bundles and connections. First, consider complex vector spaces \(V,W,U\) (everything works identically for real vector spaces). Consider the space of multilinear maps 
%     \[
%     V^*\times \overset{k}{\cdots }\times V^* \times W \overset{\ell}{\cdots} \times W \longrightarrow U,
%     \]
%     which we will denote \(\mathrm{Mult}(\times^k V^* \times ^\ell W \to U)\). There is a natural isomorphism between this set and the set of multilinear maps 
%     \[
%     V^*\times \overset{k}{\cdots }\times V^* \times W \overset{\ell}{\cdots} \times W \times U^*\to \C.
%     \]
%     The latter is naturally isomorphic to \(\otimes^k V \otimes (\otimes^\ell W^*) \otimes U\) (this is a standard result concerning the tensor product, a proof may be found in \parencite[Prop. 12.30]{lee_introSM}). Moreover, there is a natural isomorphism \(\Phi\colon \mathrm{Mult}(\times^k V^* \times ^\ell W \to U) \to \otimes^k V \otimes (\otimes^l W^*) \otimes U^*, A \to \Phi(A)\), which is given by
%     \[
%     \begin{array}{cccc}
%         \Phi(A)\colon &  V^*\times \overset{k}{\cdots }\times V^* \times W \overset{\ell}{\cdots} \times W \times U^* & \longrightarrow & \C \\
%         & (\alpha_1, \ldots \alpha_k, w_1, \ldots w_\ell, \beta) & \mapsto & \beta\big(A(\alpha_1, \ldots \alpha_k, w_1, \ldots w_\ell)\big).
%     \end{array}
%     \]
%     Thus, if \(U,W,V\) where vector bundles over a smooth manifold \(M\), the fibers of the bundle \(\otimes^k V \otimes (\otimes^l W^*) \otimes U\) can be identified with multilinear maps \( V_p^*\times \overset{k}{\cdots }\times V_p^* \times W_p \overset{\ell}{\cdots} \times W_p \times U_p^*\to \C.\) This can of course be generalized to include more vector spaces. \\

%     On the other hand, we have the usual Tensor Characterization Lemma (see \parencite[Lemma 12.24]{lee_introSM} for a proof), which states that, for a smooth manifold \(M\), a map \(\mathcal{A}\colon \mathfrak{X}(M) \times \overset{k}{\cdots} \times \mathfrak{X}(M) \to \mathscr{C}^\infty(M;\R)\) is induced by a tensor field on \(M\) if and only if it is \(\mathscr{C}^\infty(M;\R)\)-multilinear. This can be easily generalized to arbitrary vector bundles: \\
    
%     Let \(E_1, \ldots E_s, E\) be complex vector bundles over a smooth manifold \(M\). Then, a map
%     \[
%     \Gamma(E_1) \times \overset{k_1}{\cdots} \times \Gamma(E_1) \times \Gamma(E_1^*) \times \overset{\ell_1}{\cdots } \times \Gamma(E_1^*) \times \cdots \times \Gamma(E_s) \times \overset{k_s}{\cdots} \times \Gamma(E_s) \times \Gamma(E_s^*) \times \overset{\ell_s}{\cdots } \times \Gamma(E_s^*) \to \Gamma(E)
%     \]
%     is induced by a section of \(\otimes^{k_1} E_1^* \otimes (\otimes^{\ell_1} E_1) \otimes \cdots \otimes (\otimes^{k_s} E_s^*)\otimes (\otimes^{\ell_s }E_s) \otimes E^*\) if and only if it is \(\mathscr{C}^\infty(M;\C)\)-linear with respect to all of its entries. \\

%     Indeed, (suppose \(s=1\) to simplify the notation), given a section \(A \in \Gamma(\otimes^{k} E_1 \otimes (\otimes^{\ell} E_1^*) \otimes E^*)\), we can define such a map by \(\mathcal{A}(\sigma_1, \ldots \sigma_k, \omega_1 \ldots \omega_\ell)(p):=A_p(\sigma_1(p), \ldots \sigma_k(p), \omega_1(p), \ldots \omega_l(p)) \in E_p\) (here we use the identification with the space of multilinear maps). Reciprocally, if we are given \(\mathcal{A}\), we may define the section by
%     \[
%     A_p(s_1, \ldots s_k, w_1, \ldots w_\ell):=\mathcal{A}(\sigma_1, \ldots \sigma_k , \omega_1, \ldots \omega_\ell)(p),
%     \]
%     where \(\sigma_i\) and \(\omega_j\) are global smooth extensions of \(s_i \in (E_1)_p\) and \(w_j\in (E_1^*)_p\). This can always be done with a bump-function argument. \\

%     Not, back to the curvature of a connection. According to our definitions, \(R^\nabla\) is a \(\mathscr{C}^\infty(M;\C)\)-multilinear smooth map \(\Gamma(T_\C M) \times \Gamma(T_\C M) \times \Gamma(E) \to \Gamma(E)\), so by the construction above, it can be seen as a section of \(T_\C^*M \otimes T_\C^*M \otimes E \otimes E^* \). Moreover, it is antisymmetric with respect to its vector field entries, and \(E \otimes E^* \cong \mathrm{End}(E)\). Therefore, it can be seen as a section of \(\bigwedge^2 T_\C^*M \otimes \mathrm{End}(E)\), that is, as an \(\mathrm{End}(E)\)-valued \(2\)-form. When we see it as an element of \(\Omega_\C^2(M;\mathrm{End}(E))\), we will use the notation \(F_\nabla\).
% \end{obs}

% As we did with the connection, it is useful to know its coordinate expression. Let \((s_1, \ldots s_r)\) be a smooth local frame, and let \(\omega^i_j\) be the connection \(1\)-form. Again, to compute the curvature, it suffices to compute its action on the sections of the frame:
% \begin{align*}
%     R^\nabla(X,Y)(s_j) & = \nabla_X \nabla_Y s_j-\nabla_Y\nabla_X s_j -\nabla_{[X,Y]} s_j = \nabla_X(\omega^k_j(Y)s_k)-\nabla_Y (\omega^k_j(X)s_k) - \omega^k_j([X,Y])s_k \\
%     & = X(\omega^k_j(Y)) \cdot s_k + \omega^k_j(Y)\cdot \nabla_X s_k - Y(\omega^k_j(X))\cdot s_k -\omega^k_j(X) \cdot\nabla_Ys_k -\omega^k_j([X,Y])s_k \\
%     & = \big(X(\omega^k_j(Y)) - Y(\omega^k_j(X))-\omega^k_j([X,Y])\big) \cdot s_k + \omega^k_j(Y)\cdot \omega^i_k(X)\cdot s_i-\omega^k_j(X) \cdot \omega_k^i(Y) \cdot s_i \\
%     & \overset{{\color{red}(1)}}{=} \big(\diff \omega^k_j(X,Y) + (\omega^k_\ell\wedge \omega^\ell_j)(X,Y)\big) \cdot s_k 
% \end{align*}
% Where we used in \({\color{red}(1)}\) (see \parencite[Prop. 14.28]{lee_introSM})  \(\diff \omega (X,Y)=X(\omega(Y))-Y(\omega(X))-\omega([X,Y]).\)

% This motivates the definition of the \textit{curvature 2-forms}\index{Curvature!2-forms}, which are given by
% \begin{equation} \label{def_curv2_forms}
% \Omega^k_j:=\diff \omega^k_j+ \omega^k_\ell\wedge \omega^\ell_j.
% \end{equation}
% Thus, the curvature is represented locally by a matrix of \(2\)-forms, \(\Omega=(\Omega_i^j)\). In that case, we have \(\Omega=\diff \omega + \omega \wedge \omega\). Observe that \(\omega \wedge \omega\) is a wedge product of matrices, where the usual product is replaced by a wedge product (because the entries are one forms). Therefore, in general we will not have \(\omega \wedge \omega = 0\). Observe that, if \((t^1, \ldots t^r)\) is the dual reference to \((s_1, \ldots s_r)\), then we may write
% \begin{equation} \label{F_nabla_curv}
%     F_\nabla = \Omega^k_j \otimes s_k \otimes t^j.
% \end{equation}
% This follows from the fact that \(F^\nabla(X,Y,s_\ell)=R^\nabla(X,Y)(s_\ell)\).
% What is more, if we consider \(\omega^j_k\otimes s_j \otimes t^k \in \Omega_\C^1(\mathrm{End(E)})\), then \ref{def_curv2_forms} can be interpreted as an equality of \(\mathrm{End}(E)\)-valued \(2\) with the wedge product defined in {\color{red} referencia}. \\

% There is a final approach to defining the curvature, which is the one used in \parencite{kobayashi_complex_vector_bundles} and \parencite{huybrechts2005complex}. This involves considering a natural extension of the connection to a map between \(E\)-valued \(k\) and \(k+1\) forms. This is done as follows:

\Theorem{\textbf{The Exterior Covariant Derivative}}{Let \(M\) be a complex manifold, and \(\pi\colon E \to M\) be a smooth complex vector bundle over \(M\). Let \(\nabla\) be a connection on \(E\). Then, there are unique operators called \textit{exterior covariant derivatives}\index{Covariant derivatives!exterior} \(\diff_\nabla \colon \Omega_\C^q(M;E) \to \Omega_{\C}^{q+1}(M;E)\) for each \(q \geq 0\), satisfying: \label{Th_ext_cov_der}
\begin{enumerate}
    \item For every \(\sigma \in \Gamma(E)=\Omega_\C^0(M;E)\), \(\diff_\nabla \sigma=\nabla \sigma\).
    \item For every \(\alpha \in \Omega_\C^q(M)\) and \(\beta \in \Omega_\C^p(M;E)\), 
    \[
    \diff_\nabla (\alpha \wedge \beta)= \diff \alpha \wedge \beta + (-1)^q \alpha \wedge \diff_\nabla \beta.
    \]
    \item For every \(\alpha \in \Omega_\C^p(M;E)\), \(\diff_\nabla^2 \alpha = F_\nabla \wedge \alpha \in \Omega_\C^{p+2}(M;E)\).
\end{enumerate}
}
\begin{proof}
    First, observe that, with some abuse of notation, we denote all the exterior covariant derivatives the same (instead of something more reigorous like \(\diff_\nabla^q\)) to keep notation concise. \\

    Now, define, for every \(\sigma \in \Gamma(E)\), \(\diff_\nabla\sigma=\nabla \sigma\). For higher order \(E\)-valued forms, we are going to give an expression that depends on the chosen local frame, and then check that the expression is the same for any other local frame (this is similar to the way one constructs the exterior differential). Hence, let \((s_1, \ldots s_r)\) be a local reference, and let \(\alpha \in \Omega_\C^q(M;E)\) whose local expression is given by \(\alpha = \alpha ^j \otimes s_j\). We define \(\diff_\nabla \alpha:=\diff \alpha^j \otimes s_j + (-1)^q \alpha ^j \wedge \diff_\nabla s_j \). To check that this expression works for any local frame, let \((\tilde{s}_1, \ldots \tilde{s}_r)\) be another local frame. In the intersection of the corresponding open sets, by proposition \ref{prop_change_of_coordintes_frames} \(\tilde{s}_k=\tau^j_k s_j\) for some smooth maps \(\tau^j_k\), such that \(\tau=(\tau^j_k) \in \mathrm{GL}(n,\C)\). Let \(\tilde{\alpha}^j:=(\tau^{-1})_i^j \cdot  \alpha^i\). Then, we may write
    \[
    s=\alpha^j \otimes s_j = \alpha^j \otimes \parent{(\tau^{-1})^k_j \cdot \tilde{s}_k} = \parent{(\tau^{-1})^k_j \cdot \alpha^j} \otimes \tilde{s}_k =\tilde{\alpha}^k \otimes \tilde{s}_k.
    \]  
    Thus,
    \begin{align*}
        \diff_\nabla \alpha & = \diff \alpha^j \otimes s_j + (-1)^q \cdot  \alpha^j \wedge \nabla s_j \\
        & = \diff (\tau^j_k \cdot \tilde{\alpha}^k) \otimes \parent{(\tau^{-1})^\ell_j \cdot \tilde{s}_\ell} + (-1)^q \cdot (\tau^j_k \cdot \tilde{\alpha}^k) \wedge \nabla ((\tau^{-1})^\ell_j \cdot \tilde{s}_\ell) \\
        & =(\diff \tau^j_k \wedge \tilde{\alpha}^k) \otimes((\tau^{-1})^\ell_j \cdot \tilde{s}_\ell)+ ({\color{orange}\tau^j_k} \wedge \diff \tilde{\alpha}^k) \otimes ({\color{orange}(\tau^{-1})^\ell_j} \cdot \tilde{s}_\ell) \\
        & \ \ \ \ \  + (-1)^q \cdot ({\color{ForestGreen}\tau^j_k} \cdot \tilde{\alpha}^k) \wedge \big[ \diff(\tau^{-1})^\ell_j \cdot \tilde{s}_\ell + {\color{ForestGreen}(\tau^{-1})^\ell_j}  \cdot \nabla \tilde{s}_\ell\big] \\
        & = (\diff \tau^j_k \wedge \tilde{\alpha}^k) \otimes((\tau^{-1})^\ell_j \cdot \tilde{s}_\ell)+  \diff \tilde{\alpha}^k \otimes  {\color{orange}\delta_k^\ell} \cdot  \tilde{s}_\ell + (-1)^q \cdot \big[  (\tau^j_k \cdot \tilde{\alpha}^k) \wedge \diff(\tau^{-1})^\ell_j \cdot \tilde{s}_\ell + \tilde{\alpha}^k \wedge {\color{ForestGreen}\delta^\ell_k} \cdot \nabla \tilde{s}_\ell\big] \\
        & \overset{{\color{red}(2)}}{=} \diff \tilde{\alpha}^k \otimes \tilde{s}_k + (-1)^q \cdot \tilde{\alpha}^k \otimes \nabla \tilde{s}_k.      
    \end{align*}
    In \({\color{red}(2)}\), we used  
    \begin{align*}
        &(\diff \tau^j_k \wedge \tilde{\alpha}^k) \otimes((\tau^{-1})^\ell_j \cdot \tilde{s}_\ell)+ (-1)^q \cdot (\tau^j_k \cdot \tilde{\alpha}^k) \wedge \diff(\tau^{-1})^\ell_j \cdot \tilde{s}_\ell \\
        &= \big[\diff \tau^j_k \cdot (\tau^{-1})^\ell_j \wedge \tilde{\alpha}^k + (-1)^q \cdot \tau^j_k \cdot (\tilde{\alpha}^k \wedge \diff(\tau^{-1})^\ell_j ) \big] \otimes \tilde{s}_\ell \\
        &= \big[ \diff \tau^j_k \cdot (\tau^{-1})^\ell_j+(-1)^{2q+1} \cdot \tau^j_k \cdot \diff(\tau^{-1})^\ell_j \big] \wedge \tilde{\alpha}^k \otimes \tilde{s}_\ell = \diff(\tau^j_k \cdot (\tau^{-1})^\ell_j) \wedge \tilde{\alpha}^k \otimes \tilde{s}_\ell \\& = \diff \parent{\delta^\ell_k\wedge \tilde{\alpha}^k} \otimes \tilde{s}_\ell=0.
    \end{align*}
    Thus, \(\diff_\nabla\) is well-defined, and by definition, 1 trivially is satisfied. To check 2, let \(\beta\in \Omega_\C^p(M;E)\) whose expression in the domain of a local frame is \(\beta=\beta^j\otimes s_j\), and let \(\alpha \in \Omega^q_\C(M)\). Then, 
    \begin{align*}
        \diff_\nabla(\alpha \wedge \beta) & = \diff_\nabla((\alpha\wedge \beta^j) \otimes s_j) = \diff(\alpha \wedge \beta^j) \otimes s_j+(-1)^{p+q} \cdot (\alpha \wedge \beta^j) \otimes \nabla s_j \\
        & =\big[ \diff \alpha \wedge \beta^j + (-1)^q \cdot \alpha \wedge\diff  \beta^j \big] \otimes s_j + (-1)^{p+q} \cdot (\alpha \wedge \beta^j) \otimes\nabla s_j \\
        & = \diff \alpha \wedge (\beta^j \otimes s_j) + (-1)^q \cdot\alpha \wedge  \big[ \diff \beta^j \otimes s_j + (-1)^p \cdot \beta^j \otimes \nabla s_j\big] \\
        & = \diff \alpha \wedge \beta + (-1)^q \cdot \alpha \wedge \diff_\nabla \beta.
    \end{align*}
    Finally, to check 3, let \(\alpha \in \Omega_\C^p(M;E)\), and \(\sigma \in \Gamma(E)\). By 2, we know that (observe that \(\diff_\nabla\) s linear by definition)
    \begin{align*}
        \diff_\nabla^2 (\alpha \wedge \sigma) & =\diff_\nabla \parent{ \diff \alpha \wedge \sigma + (-1)^p \cdot \alpha \wedge \nabla \sigma}= \diff_\nabla(\diff \alpha \wedge \sigma) + (-1)^p \cdot \diff_\nabla (\alpha \wedge \nabla\sigma) = \\
        &= \diff^2 \alpha \wedge \sigma -(-1)^p \cdot  \diff \alpha \wedge \nabla \sigma+(-1)^p \cdot \big[ \diff \alpha \wedge \nabla \sigma +(-1)^p \cdot  \alpha \wedge \diff_\nabla^2 \nabla\big] \\
        & =  \alpha \wedge \diff_\nabla^2 \sigma.
    \end{align*}
    There are two conclusions we may extract from the previous calculation. The first one is that the action of \(\diff_\nabla^2\) on all the spaces \(\Omega_\C^p(M;E)\) is determined by its action on \(\Gamma(E)\) (every \(E\)-valued \(p\)-form is a sum of elements of the form \(\alpha \otimes \sigma=\alpha \wedge \sigma\)). The second one is that it suffices to check 3 for the elements of a local frame. Indeed, let \((s_1, \ldots s_r)\) a local frame. Then, if we assume that 3 holds for the elements of the frame, for every \(\alpha=\alpha^j \wedge s_j=\alpha^j \otimes s_j\), we get
    \[
    \diff_\nabla^2(\alpha) = \diff_\nabla^2(\alpha^j \wedge s_j) = \alpha^j \wedge \diff_\nabla^2 s_j =\alpha^j \wedge (F_\nabla \wedge s_j)= F_\nabla \wedge \alpha,
    \]
    where the last identity follows from checking that the coordinate expression of both sides is equal to \((\Omega^\ell_j \wedge \alpha^j)\otimes s_\ell\). \\

    Finally, let us check that 3 holds for the sections of a local frame:
    \begin{align*}
        \diff_\nabla^2(s_j) & = \diff_\nabla(\nabla s_j)=\diff_\nabla (\omega^k_j \otimes s_k) =\diff \omega^k_j \otimes s_k + \omega ^k_j \wedge \nabla s_k=(\diff \omega^k_j+\omega_j^\ell \wedge \omega^k_\ell) \otimes s_k \\
        & = \Omega^k_j \otimes s_k = \Omega_j^k \wedge s_k.
    \end{align*}
    Finally, from equation \eqref{F_nabla_curv} we get
    \[
    F_\nabla \wedge s_j = (\Omega^k_\ell \otimes s_k \otimes t^\ell) \wedge s_j = \delta^\ell_j \cdot \Omega^k_l \otimes s_k = \Omega^k_j \otimes s_k.
    \]
\end{proof}

Some sources then define the curvature of the connection to be the map \(\diff_\nabla^2 \colon \Gamma(E) \to \Omega_\C^2(M;E)\). It is common to find the notation \(F_\nabla=\nabla \circ \nabla\), which is a slight abuse of notation that comes from using the symbol \(\nabla\) to refer to all maps \(\diff_\nabla\). Observe that for sections \(\xi \in \Gamma(M;E)\), our definition of the wedge product means that we can write \(\diff_\nabla^2(\xi)=F_\nabla(\xi)\). \\

% \begin{obs}
%     Observe that, in terms of coordinates, we have \(\diff_\nabla^2 s=s \otimes \Omega\). Moreover, if we choose another local frame \(\tilde{s}=s \cdot \tau\), then \(\Omega'=\tau^{-1} \Omega \tau\).
% \end{obs}

\section{Induced connections} \label{sec_induced_connections}
Connections on vector bundles give rise to connections in other vector bundles, provided these are constructed using the ones we originally had. We will be concerned with 3 cases:
\begin{enumerate}
    \item Let \(E\) and \(F\) be two vector bundles over a smoth manifold, endowed with connections \(\nabla^E\) and \(\nabla^F\). Then, we may define a connection in \(E \times F\) by extending by linearity the following expression
    \begin{equation} \label{def_connection_product}
    \parent{\nabla^E \otimes \nabla^F}(\xi \otimes \eta)=\nabla^E \xi \otimes \eta + \xi \otimes \nabla^F \eta,
    \end{equation}
    as a direct computation shows.
    \item Let \(\nabla\) be a connection on \(E\), then, it induces a connection on \(E^*\) as follows. There is a unique connection \(\nabla^*\) on \(E^*\) such that for every \(t \in \Gamma(E^*)\), and \(\xi \in \Gamma(E)\), 
    \begin{equation} \label{def_connection_dual}
        \big[\nabla^*(t) \big](\xi) = \diff (t (\xi)) - t(\nabla \xi),
    \end{equation}
    where \(t(\alpha \otimes s):=t(s) \cdot \alpha\). Combining this with the previous item, we get an induced connection in \(\mathrm{Hom}(E, F)\cong E^*\otimes F\).
    \item A connection \(\nabla\) on a complex vector bundle \(E\) induces a connection on \(\bar{E}\) given by \(\overline{\nabla} s:=\overline{\nabla \overline{s}}\). The map \(E \to \bar{E}, s \mapsto \bar{s}\) sends every element to itself, but it is antilinear: \(\overline{\lambda *s}=\overline{\overline{\lambda}\cdot s}=\bar{\lambda} \cdot s\).
\end{enumerate}

% \begin{obs}
%     Equation \ref{def_connection_dual} does not define the dual connection completely on its own. Indeed, for \(\nabla^*t\) to be a section of the bundle \(T_\C^*M \otimes E^*\), we have to say how \((\nabla^*t)_p\) acts on a vector \(v \in E_p\) to give an element \((T_pM)\otimes _\R \C\). This is done by defining \((\nabla^*t)_p(v):=[(\nabla^*t)(\xi)]_p\) for any smooth section \(\xi\) such that \(\xi_p=v\). However, we have to check that this does not depend on the choice of smooth section. One way to this is using local frames. Clearly, it suffices to check that for a smooth section \(\xi\) satisfying \(\xi_p=0\), we have \([(\nabla^*t)(\xi)]_p=0\) Let \((s_1, \ldots s_r)\) be a local frame, and \((t^1, \ldots t^r)\) be its dual frame. For some smooth complex-valued coordinate functions, we can write \(\xi=\xi^i  \cdot s_i\) (note that \(\xi^i(p)=0\) for all \(i\)) and \(t=\lambda_j \cdot t^j\). Then,
%     \begin{align*}
%         [(\nabla^*t)(\xi)]_p&=\left. \diff(t(\xi)) - t(\nabla \xi) \  \right|_{p} = \diff (\lambda_j\cdot \xi^j)-t(\diff \xi^j\cdot s_j+\xi^j \otimes s_j) \ |_{p} \\
%         & = \lambda_j(p) \cdot \diff_p \xi^j+\xi^j(p) \cdot \diff_p \lambda^j - \diff_p\xi^j \cdot t(s_j)-\xi^j(p)\cdot t(s_j) =0.
%     \end{align*}
% \end{obs}

\section{Hermitian inner products}
Hermitian inner products can be seen as generalizations of Riemannian metrics to arbitrary vector bundles.

\Definition{\textbf{Hermitian fiber metric}}{Let \(\pi \colon E \to M\) a smooth complex vector bundle over a smooth manifold \(M\). A \textit{hermitian fiber metric}\index{Hermitian!fiber metric} \(h\) is an assignment of a hermitian inner product \(h_x\) to each fiber \(E_x\) and such that for every pair of local smooth sections \(\sigma, \tau \in \Gamma(U;E)\), the map \(U \to \C, x \mapsto h_x(\sigma_x, \tau_x)\) is smooth. A vector bundle that is endowed with a hermitian fiber metric is called a \textit{hermitian vector bundle}\index{Hermitian!vector bundle} \label{def_hermitian_fiber_metric}}
That is, a hermitian fiber metric is just an assignment of a hermitian inner product to each fiber, with the condition that is varies smoothly. Under some abuse of language, we will often call \(h\) an \textit{hermitian inner product}\index{Hermitian!inner product}. We adopt the convention that hermitian inner products are antilinear with respect to their second entry, i.e. \(h(v, \lambda \cdot u)= \bar{\lambda} \cdot h(v,u)\)for \(\lambda \in \C\). Of course, there is analogous idefinition for real vector bundles, where we assign an inner product to each fiber that varies smoothly. Riemannian metrics are examples of inner products in the bundle \(TM\). 

% \begin{obs}
%     If \(h\) is a Hermitian inner product on a complex vector bundle \(E\), we may see it as a section of the vector bundle \((E \otimes \bar{E})^*\). Indeed, observe that for each \(x \in M\), \(h_x\) can be seen as a \textit{bilinear} map \(h\colon E_x \times \bar{E}_x \to \C\) (sesguilinearity is now encoded in \(\bar{E}_x\)). To see this, just define
%     \[
%     E_x \times \bar{E}_x \to \C, (u,v) \mapsto h_x(u,v)
%     \] 
%     and observe that \(h_x(u, \lambda*v)=h_x(u, \overline{\lambda} \cdot v)=\lambda \cdot h(u,v)\). The set of all bilinear maps \(E_x \times \bar{E}_x \to \C\) defines a \(\C\)-vector space that is isomorphic to \((E_x \otimes \bar{E}_x)^*\). We thus have a rough section of the bundle \((E \otimes \bar{E})^*\) which is smooth thanks to the condition of the definition.
% \end{obs}

If \(E\) is endowed with a hermitian fiber metric \(h\), we say that a connection is \textit{compatible with the fiber metric}\index{Connection!compatible with the fiber metric}, \textit{compatible with the hermitian inner product} or \textit{unitary}\index{Connection!unitary} if it satisfies the following relation (\(U \subset M\) is some open set)
\[
X \big(h(\xi, \eta) \big)=h(\nabla_X \xi, \eta)+h(\xi , \nabla_{\bar{X}} \eta) \hspace{0.75 cm} \text{ for all } \hspace{0.75 cm} X \in \Gamma(U, T_\C M), \ \ \xi, \eta \in \Gamma(U;E).
\]
This is equivalent to the following condition
\begin{equation} \label{def_cond_compatibility}
\diff h(\xi,\eta)=h(\nabla \xi, \eta) + h(\xi, \nabla \eta).
\end{equation}
Let us comment on the terms of equation \eqref{def_cond_compatibility} above. For any pair of smooth sections \(\xi, \eta \in \Gamma(U;E)\), the definition of hermitian fiber metric implies that \(h(\xi,\eta)\) is smooth, hence \(\diff h(\xi, \eta)\) is a smooth one form. On the right-hand side, we have \(h(\nabla \xi, \eta)\) (and a symmetric term). Now, \(\nabla \xi\) is a section of the bundle \(T^*_\C M \otimes E\). The product of a section of \(T^*M \otimes E\) and one of \(E\) through \(h\) is, by definition \(h(\alpha \otimes \xi, \eta):=h(\xi, \eta) \cdot\alpha\), \(h(\xi, \alpha \otimes \eta):=h(\xi,\eta)\cdot\bar{\alpha}\) (this is the conjugation of forms defined in section \ref{sec_Dolbeault_complex}), and it is extended by sesguilinearity to general sections of \(T^*_\C M \otimes E\).

% \begin{obs}
%     A global section \(h\) of the bundle \((E\otimes \bar{E})^*\) induces a hermitian inner product if and only if \(\nabla h=0\), where by \(\nabla\) we also mean the natural connection induced by the original connection on \((E \otimes \bar{E})^*\). Indeed, a direct calculation shows that \(\nabla h(\xi, \eta)=\diff h(\xi, \eta)-h(\nabla \xi, \eta)-h(\xi,\nabla \eta)\).
% \end{obs}

\section{The First Real Chern Class}
\Definition{\textbf{First Real Chern Class}}{Let \(\pi\colon E \to M\) be a complex vector bundle over a smooth manifold \(M\), and let \(\nabla\) be a connection. We define the \textit{First Chern form}\index{Chern!form} of \(\nabla\) as 
\[
c_1(\nabla):=\frac{\ii}{2 \pi}\  \mathrm{tr} \ F_\nabla.
\]}
In the previous definition, the trace of an element \(\alpha \otimes f \in \Omega_\C^2(M; \mathrm{End}(E))\) is defined as \(\mathrm{tr}(\alpha \otimes f):=\mathrm{tr}(f) \cdot \alpha \in \Omega_\C^2(M)\).

\Theorem{}{Let \(\pi\colon E \to M\) be a complex vector bundle. Then, \(c_1(\nabla)\) is closed, and the de Rham cohomology class of \(c_1(\nabla)\) in \(H^2_{\mathrm{dR}}(M;\C)\) is independent of the choice of the connection.}
\begin{proof}
    Choose a connection \(\nabla\) (they always exist by equation \eqref{def_nontrivial_connection}), a local frame \((s_1, \ldots s_r)\) and its dual frame \((t^1, \ldots t^r)\). Since \(F_\nabla=\Omega_j^i\otimes s_i \otimes t^j\), we have \(\mathrm{tr} (F_\nabla)=\Omega^i_j \cdot \delta^j_i=\Omega^j_j\) (sum implied). Thus, 
    \[
    c_1(\nabla)=\frac{\ii}{2\pi} \cdot \Omega^j_j=\frac{\ii}{2\pi} \cdot \parent{\diff \omega^j_j+\omega^j_\ell \wedge \omega^\ell_j}=\frac{\ii}{2\pi} \cdot \diff \omega^j_j,
    \]
    where the second summand is zero because \(\omega^j_\ell \wedge \omega^\ell_j=-\omega_j^\ell\wedge \omega_\ell^j=-\omega_\ell^j \wedge\omega_j^\ell\) (we renamed the dummy indexes in the last step). Hence, \(c_1(\nabla)\) is locally exact, and therefore globally closed (recall that \(\diff\) acts locally). It need not be globally exact though, because the connection one forms are defined only locally. To check all connections define the same cohomology class, let \(\nabla'\) be another connection. Define the map \(\mathscr{D}\colon \Gamma(T_\C M) \times \Gamma(E) \to \Gamma(E), (X, \sigma) \mapsto \nabla_X \sigma-\nabla'_X \sigma\). It is \(\mathscr{C}^\infty(M;\C)\)-linear with respect to all of its components (see equation \eqref{eq_D_linear} below, for the \(X\)-componet it is obvious), so it defines an element of \(\Omega^1_\C(M; \mathrm{End}(E))\). 
    \begin{equation} \label{eq_D_linear}
        \mathscr{D}_X(f \sigma)=\nabla_X (f \cdot \sigma)-\nabla'_X(f \cdot \sigma)=\diff f \otimes \sigma + f \cdot \nabla_X \sigma-\diff f \otimes \sigma-f \cdot \nabla'_X \sigma= f \cdot \mathscr{D}_X(\sigma)
    \end{equation}
    For every local frame, we have \(\mathscr{D}=(\omega^j_i-(\omega')^j_i) \otimes s_j \otimes t^i\), hence
    \begin{align*}
        c_1(\nabla)-c_1(\nabla')=\frac{\ii}{2 \pi} \parent{\diff \omega^j_j-\diff (\omega')^j_j}=\diff \parent{\frac{\ii}{2 \pi} \cdot \mathrm{tr} \mathscr{D}},
    \end{align*} 
    because the coordinate expression we used is valid for every local frame. Therefore, \(c_1(\nabla)\) and \(c_1(\nabla')\) differ by an exact \(2\)-form, as we wanted to see.
\end{proof}

The following proposition will be needed to define the first Chern class, and it is what explains the factor \(\ii\) in the definition of the Chern form.
\Proposition{}{Let \(\pi\colon E \to M\) be a complex vector bundle over a smooth manifold which is endowed with a hermitian fiber metric. If \(\nabla\) is a metric connection on \(E\), then \(c_1(\nabla)\) is a real \(2\)-form.}
\begin{proof}
    Applying the Gram-Schmidt algorithm (see \parencite[Lemma 8.3]{lee_introSM}) to a frame around any point, we may always obtain an orthonormal frame, i.e. \(h(s_i, s_j)=\delta_{ij}\). If \(\nabla\) is a metric connection, then \(\omega^k_j=-\overline{\omega^j_k}\). Indeed, for every local vector field, the compatibility condition implies (no summation in the last expression)
    \begin{align*}
        0 & = X(h(s_i, s_j)) = h(\nabla_X s_i, s_j) + h(s_i,\nabla_X s_j) = h(\omega^k_i(Z)s_k, s_j)+h(s_i, \omega^k_j(X)s_k) = \omega^j_i(X)+\overline{\omega^i_j(X)}. 
    \end{align*}
    We thus find \(\omega^k_j+\overline{\omega^j_k}=0\) (no summation). Now, considering \(j=k\) and taking the sum over \(j\), \(\omega^j_j+\overline{\omega^j_j}=0\). Hence, \(\omega^j_j\) is purely imaginary, thus, so is \(\Omega^j_j\), and therefore, \(c_1(\nabla)=\frac{\ii}{2 \pi} \cdot \Omega_j^j\) is real.
\end{proof}

Observe that, in general, \(c_1(\nabla)\) defines an element of \(H^2_{\mathrm{dR}}(M;\C)\). However, the previous proposition shows that whenever \(\nabla\) is compatible with the metric, \(c_1(\nabla)\) is real. Since we may always find a connection compatible with metric (for instance, the Chern connection, see theorem \ref{th_Chern_connection}), we conclude that \(c_1(\nabla)\) defines (for \textit{any connection}) an element of the real De Rham cohomology group \(H^2_{\mathrm{dR}}(M;\R)\) (which we may see as the subspace of \(H^2_{\mathrm{dR}}(M;\C)\) of cohomology classes invariant under conjugation, see section \ref{sec_Dolbeault_complex}). 


\section{The Chern connection}
Let \(\pi\colon E \to M\) be a holomorphic vector bundle over a complex manifold \(M\). Recall that in this case, there exists a Cauchy-Riemann operator \(\overline \partial_E\colon \Omega_\C^{(p,q)}(M;E) \to \Omega_\C^{(p,q+1)}(M;E)\) (see proposition \ref{Cauchy_Riemann_operator}). Since any connection \(\nabla\) can be seen as a map \(\nabla\colon \Omega_\C^0(M;E)\to \Omega_\C^1(M;E)\), and there is a direct sum decomposition, \(\Omega^1_\C(M;E)=\Omega^{(0,1)}_\C(M;E)\oplus \Omega_\C^{(1,0)}(M;E)\), we can split the connection into two components 
\[
\nabla^{(1,0)}\colon \Omega_\C^0 (M;E) \to \Omega_\C^{(1,0)}(M;E) \hspace{1.5 cm}
\nabla^{(0,1)}\colon \Omega_\C^0(M;E)\to \Omega_\C^{(0,1)}(M;E),
\]
and therefore we may write \(\nabla=\nabla^{(1,0)}+\nabla^{(0,1)}\). This is just the composition of \(\Pi^{(1,0)}\) and \(\Pi^{(0,1)}\) with \(\nabla\) (extending each projection to \(T_\C^*M \otimes E\)).

\begin{obs} \label{nabla01_is_Dolbeault_operator}
    Since \(\nabla\) is a connection, it satisfies its "Leibniz rule", and so \(\nabla(f\cdot s)=\diff  f \otimes s+f\cdot \nabla s\). Since \(\diff=\partial + \barpartial\), after decomposition, we find, for all \(f \in \mathscr{C}^\infty(M;\C)\) and \(s \in \Omega_\C^{0}(M;E)\)
    \[
    \nabla^{(0,1)}(f\cdot s)= \barpartial f \otimes s+f\cdot \nabla^{(0,1)} s.
    \]
    This is the same relation satisfied by \(\overline{\partial}_E\) (see \eqref{prop_operador_cauchy_riemann}). This motivates the following definition.
\end{obs}

\Definition{\textbf{Compatible connection}}{A connection \(\nabla\) on a holomorphic vector bundle \(\pi\colon E \to M\) is \textit{compatible with the holomorphic structure} if \(\nabla^{(0,1)}=\barpartial_E\)}

Our goal is to show that, for a holomorphic hermitian vector bundle, there is a special connection, the Chern connection, that is compatible with both, the holomorphic structure and the hermitian fiber metric. To do that, it is useful to have several equivalent characterizations of being compatible with the holomorphic structure. This is the content of the following lemma.

\Lemma{}{Let \(\pi \colon E \to M\) be a complex vector bundle over a complex manifold \(M\), and let \(\nabla\) be a connection on \(E\). The following are equivalent:
\begin{enumerate}
    \item \(\nabla\) is compatible with the holomorphic structure (\(\nabla^{(0,1)}=\barpartial_E\)).
    \item If \(U \subset M\) is open, \(\xi \in \Gamma(U;E)\) is holomorphic, and \(\bar{Z}\) is a smooth local section of \(T^{(0,1)}M\), then \(\nabla_{\bar{Z}} \xi = 0\).
    \item For each holomorphic frame, the connection forms are of type \((1,0)\).
\end{enumerate}
}
\begin{proof}
    All the statements are local, so we may just work in the domain of a holomorphic local frame \((s_1, \ldots s_r)\). Let \(\xi\) be a smooth local section defined on the same open subset \(U\). Write \(\xi=\xi^i s_i\). Then, from remark \ref{coordinate_exp_connection}, we have \(\nabla \xi=\diff \xi ^j \otimes s_j+ \xi^j \cdot \omega^k_j \otimes s_k\). Projecting onto \({\bigwedge}^{(0,1)}_\C M\otimes E\), we find,
    \[
    \Pi^{(1,0)} \parent{\nabla\xi}=\nabla^{(0,1)}\xi = \barpartial \xi^j \otimes s_j + \xi^j \cdot (\omega_j^k)^{(0,1)} \otimes s_k. 
    \]
    On the other hand, by definition of the Cauchy-Riemann operator for vector bundles, we have \(\barpartial_E \ \xi=\barpartial \xi^j \otimes s_j\). Observe that \(\barpartial_E = \nabla^{(0,1)}\) if and only if \(\xi^j \cdot (\omega^k_j)^{(0,1)} \otimes s_k=0\) for every local section \(\xi\), that is, if and only if \((\omega^k_j)^{(0,1)}=0\) for all \(j,k\). This shows that \((1)\) and \((3)\) are equivalent. Suppose now that \((3)\) holds, and let \(\xi=\xi^j\cdot s_j\) be a holomorphic local section of \(E\) (observe that the maps \(\xi^j\) are now holomorphic), and \(Z \in \Gamma(U; T^{(0,1)}M)\). Then,
    \begin{equation} \label{eq1_lemma_chern_connection}
    \nabla_{Z}\  \xi=Z(\xi^j) \cdot s_j+\xi^j \cdot \omega^k_j(Z) \cdot s_k  =0.
    \end{equation}
    Where we used the comments in remarks \ref{holomorp_func_zero} and \ref{charac_pq_forms}. Finally, suppose that 2 holds, then, equation \eqref{eq1_lemma_chern_connection} implies that \(\omega^k_j(Z)=0\) for every \(Z \in \Gamma(U, T^{(0,1)}M)\), hence  by remark \ref{charac_pq_forms}, each \(\omega^k_j\) is a \((1,0)\)-form.
\end{proof}

\Theorem{\textbf{Existence of the Chern connection}}{Let \(\pi\colon E \to M\) be a complex vector bundle over a complex manifold \(M\). There is a unique connection, called the \textit{Chern connection}\index{Chern!connection} that is compatible with the metric and the holomorphic structure. \label{th_Chern_connection}}
\begin{proof}
    The strategy is to show uniqueness first, and use what we will find as a hint on how to define the connection. Let \(\nabla\) be a connection compatible with the holomorphic structure. By the previous lemma, the connection one-forms with respect to a holomorphic frame \((s_1, \ldots s_r)\) are of type \((1,0)\). Let \(Z\) be a holomorphic section of \(T^{(1,0)}M\). Then, \(\bar{Z}\) is a smooth section of \(T^{(0,1)}M\). Let \(h_{jk}:=h(s_j, s_k)\). Compatibility with the metric \({\color{red}(3)}\) and the connection \({\color{red}(4)}\) imply
    \begin{align} \label{eq1_Chern_connection}
        Z(h_{jk}) & = Z(h(s_j, s_k)) \overset{{\color{red}(3)}}{=} h(\nabla_Z s_j, s_k) +h(s_j, \nabla_{\bar{Z}} s_k) \overset{{\color{red}(4)}}{=}h(\omega^i_j(Z)s_i, s_k)=h_{ik}\cdot \omega^i_j(Z).
    \end{align}
    This is valid for every holomorphic section \(Z\) of \(T^{(1,0)}M\). Observe that \(h_{jk}\) is a smooth function and \(\diff h_{jk}(Z)=Z(h_{jk})\). Moreover, since \(\diff=\partial+\barpartial\), and \(\barpartial h_{jk}(Z)=0\), we have \(\partial h_{jk}(Z)=h_{ik} \cdot \omega^i_j(Z)\). Since \(h\) is a hermitian fiber metric, the matrix \(h_{ik}\) has an inverse, which we denote, as usual, by \(h^{ik}\). Multiplying by that inverse on both sides, we get \(h^{km} \partial h_{jk}(Z)=h^{km}h_{ik} \cdot \omega^i_j(Z)=\delta^m_i \cdot \omega^i_j(Z)=\omega^m_j(Z)\). Finally, since this is valid for every section \(Z\) of \(T^{(1,0)}M\), we get 
    \begin{equation} \label{one_form_Chern_connection}
    \omega^m_j=h^{km} \partial h_{jk}.
    \end{equation}
    Observe that the right-hand side is a \((1,0)\)-form, which was to expect after the result of the previous lemma. We have determined the connection one forms in terms of the coefficients \(h_{jk}\) for every holomorphic local frame. Since a connection is uniquely given (at least locally) by its connection one forms, we have shown uniqueness. \\

    Now, to show that it exists, the key idea is to use equation \ref{one_form_Chern_connection} to define a connection in the domain of each local holomorphic frame. Then, uniqueness ensures that the connection is well-defined where two such frames overlap, and thus, the connection is globally defined. Therefore, we can define \(\nabla\) in terms of each holomorphic local frame, and these definitions can be glued together to produce a globally defined connection. It remains to check that it is compatible with the metric and the holomorphic structure. Compatibility with the holomorphic structure follows from the previous characterization lemma, together with the fact each \(\omega^m_j=h^{km} \partial h_{jk}\) is a \((1,0)\)-form. Finally, compatibility with the metric can be deduced by direct computation as follows. Let \((s_1, \ldots s_r)\) be an holomorphic local frame, and \(\xi = \xi^i \cdot s_i\) and \(\eta=\eta^j\cdot s_j\) smooth sections defined on the same open set. Let \(Z\) be a complex vector field, which may decomposed as \(Z=X+Y\), where \(X\) is a section of \(T^{(1,0)}M\), and \(Y\) is a section of \(T^{(0,1)}M\). Equation \ref{eq1_Chern_connection} applies to \(X\), and the analogous equation for \(Y\) is
    \begin{equation} \label{eq2_Chern_connection}
        Y(h_{ij})=h(\nabla_{ Y} s_i, s_j)+h(s_i, \nabla_{\bar{Y}} s_j)=h(s_i, \omega^k_j(Y)s_k)=\overline{\omega^k_j(Y)} s_k.
    \end{equation}
    Now, on the one hand,
    \begin{align} \label{eq3_Chern_connection}
        \begin{split}
        Z(h(\xi, \eta)) & =Z(h_{ij} \xi^i \overline{\eta^j})=\xi^i \overline{\eta^j} Z(h_{ij}) + h_{ij} Z(\xi ^i \overline{\eta^j} )=\xi^i\overline{\eta^j}\cdot \parent{X(h_{ij}) + Y(h_{ij})} + h_{ij} Z(\xi ^i \overline{\eta^j} ) \\
        & = \xi^i  \overline{\eta^j}\cdot \parent{\omega^k_j(X) + \overline{\omega^k_j(\bar{Y})}}\cdot h_{ik} + h_{ij} \cdot \parent{Z(\xi^i)\overline{\eta^j} + \xi^i Z(\overline{\eta^j})}.
        \end{split}
    \end{align}
    On the other hand,
    \begin{align} \label{eq4_Chern_connection}
        \begin{split}
        h(\nabla_Z \xi, \eta) + h(\xi, \nabla_{\bar{Z}} \eta) & = h\parent{Z(\xi^i)s_i + \xi^i \nabla_Z s_i, \eta} +h\parent{\xi, \bar{Z}(\eta^j)s_j+\eta^j \nabla_{\bar{Z}}s_j}\\
        & = Z(\xi^i) \overline{\eta^j}h_{ij}+ \xi ^i \overline{\eta^j}\cdot {\omega_i^k(X)}\cdot h_{ik} + \xi^i \cdot  \overline{\bar{Z}(\eta^j)}\cdot h_{ij} + \xi^i  \overline{\eta^j} \cdot \overline{ \omega^k_j(\bar{Y})}\cdot h_{ik} \\
        & \overset{{\color{red}(5)}}{=} \parent{Z(\xi^i) \overline{\eta^j} + \xi^i Z(\overline{\eta^j})}\cdot h_{ij} + \xi^i \overline{\eta^j } \cdot h_{ik} \cdot \parent{\omega^k_j(X) + \overline{\omega^k_j(\bar{Y})}}.
        \end{split}
    \end{align}
    In \({\color{red}(5)}\) we used that \(\overline{\bar{Z}(\eta^j)}=Z(\overline{\eta^j})\). We thus get that the last expressions in equations \eqref{eq3_Chern_connection} and \eqref{eq4_Chern_connection} coincide, hence the connection is compatible with the hermitian metric.
\end{proof}
Observe that this proof is similar to the one given to construct the Levi-Civita connection; there, one writes the Christoffel symbols in terms on the coefficients of the riemannian metric on each local chart. When we work with several connections, we will denote by \(\nabla_H\) the Chern connection. 

As we anticipated before the previous theorem allows us to define the \textit{First Real Chern class}\index{Chern!class} of \(E\) as \(c_1(E):=[c_1(\nabla)]\in H^2_{\mathrm{dR}}(M;\R)\). This means that, associated to any holomorphic vector bundle, there is a real cohomology class that \textit{does not depend} on the connection chosen, only on the vector bundle itself. \\

To close this chapter, we will show a very important property satisfied by the curvature of the Chern connection: it is of type \((1,1)\).

\Proposition{\textbf{Curvature of the Chern connection}}{Let \(M\) be a complex manifold, and \(\pi\colon E \to M\) be a complex vector bundle over \(M\) endowed with a hermitian fiber metric. Let \(\nabla\) be its Chern connection. Then, \(F_{\nabla} \in \Omega_\C^{(1,1)}(M; \mathrm{End}(E))\).}
\begin{proof}
    From equations \eqref{def_curv2_forms} and \eqref{F_nabla_curv}, we know that, with respect to a local frame, we may write
    \[
    F_\nabla = (\partial \omega^k_j + \barpartial \omega^k_j + \omega^k_\ell \wedge \omega^\ell_k) \otimes s_k \otimes t^j,
    \]
    so we just need to check that, for the Chern connection, the curvature two-forms \(\Omega_j^k=\partial \omega^k_j + \barpartial \omega^k_j + \omega^k_\ell \wedge \omega^\ell_k\) is of type \((1,1)\). We now that the connection one forms of the Chern connection are of type \((1,0)\). This means that \(\barpartial \omega^k_j \) is of type \((1,1)\), and \(\partial \omega^k_j + \omega^k_\ell \wedge \omega^\ell_k\) is of type \((2,0)\). Therefore, it suffices to show that the latter vanishes. This is a calculation using the expression of the connection one forms in terms of the metric (see equation \eqref{one_form_Chern_connection}). 
    \begin{equation} \label{eq1_F_nabla_11}
        \begin{split}
        \partial \omega^k_j + \omega^k_\ell \wedge \omega^\ell_k & = \partial (h^{km} \partial h_{jk})+ (h^{pm} \partial h_{kp}) \wedge (h^{qk} \partial h_{jq})  \\
        & = \partial h^{km} \wedge \partial h_{jk} + h^{pm}h^{qk} \cdot \partial h_{kp} \wedge \partial h_{jq}.
    \end{split}
    \end{equation}
    On the other hand, 
    \[
    0=\partial\parent{\delta ^k_\ell} = \partial \parent{h^{kp}h_{p\ell}}=\partial h^{kp} \cdot h_{p \ell} + h^{kp} \cdot \partial h_{p \ell},
    \]
    therefore, \(h_{p \ell} \cdot \partial h^{hp}=-h^{kp} \cdot \partial h_{p \ell}\). Multiplying on both sides by \(h^{\ell m}\) and taking the sum yields, 
    \[
    h^{\ell m} h_{p \ell} \cdot \partial h^{kp}=\delta^m_p \partial h^{kp}\overset{{\color{red}(6)}}{=}\partial h^{km}=-h^{\ell m} h^{kp} \cdot \partial h_{p\ell}.
    \]
    Taking equality \({\color{red}(6)}\) and inserting it in equation \eqref{eq1_F_nabla_11}, yields
    \begin{align*}
        \partial h^{km} \wedge \partial h_{jk} + h^{pm}h^{qk} \cdot \partial h_{kp} \wedge \partial h_{jq}= -h^{\ell m}h^{kp} \cdot \partial h_{p \ell} \wedge \partial h_{jk}+ h^{pm}h^{qk} \cdot \partial h_{kp} \wedge \partial h_{jq}=0.
    \end{align*}
\end{proof}

\section{The Chern correspondence} \label{sec_Chern_correspond}
We devote this section to explaining, schematically, the Chern correspondence. Good references for this section are \parencite[Ch. 7]{kobayashi_complex_vector_bundles} and \parencite[Ch. 1]{teleman}. The main statement of the Chern correspondence is that there is a bijection between the set of connections on a complex vector bundle over a complex manifold and the set of Dolbeault operators (or semiconnections), which restricts to a bijection between connections of type \((1,1)\), and integrable Dolbeault operators. 

\Definition{\textbf{Dolbeault operator}}{Let \(E \to M\) be a complex vector bundle over a complex manifold. A \textit{Dolbeault operator}\index{Dolbeault!operator} or \textit{semiconnection of type \((0,1)\)} on \(E\) is a \(\C\)-linear map
\[
\barpartial_e \colon \Gamma(M; E) \longrightarrow \Omega_\C^{(0,1)}(M;E)
\]
satisfying
\begin{equation} \label{eq_ext_barpartial_e}
\barpartial_e(f \cdot \xi) = \barpartial f \otimes \xi + f \cdot \barpartial_e \xi.
\end{equation}
We denote \(\mathscr{D}(E)\) (or \(\mathscr{D}\) when the bundle is clear from the context) the set of all Dolbeault operators.
}
Observe that every Dolbeault operator can be extended to higher order complex differential forms by linear extension of
\begin{equation} \label{eq_ext_Dolb_operator}
\barpartial_e (\alpha \otimes \xi) = \barpartial \alpha \otimes \xi + (-1)^{p+q} \cdot \alpha \wedge \barpartial_e \xi, \hspace{1 cm} \alpha \in \Omega_\C^{(p,q)}(M), \ \xi \in \Gamma(M; E). 
\end{equation}

\Lemma{}{Let \(E \to M\) be a complex vector bundle over a complex manifold. Then, \(\mathscr{D}(E)\) is an affine space modelled on the vector space \(\Omega_\C^{(0,1)}(\mathrm{End}(E))\). \label{lemma_vect_bund_connections_affine_space}}
\begin{proof}
    First, fix \(\barpartial_{e}^0 \in \mathscr{D}(E)\), and choose \(\alpha \in \Omega_\C^{(0,1)}(\mathrm{End}(E))\). Let us see that \(\barpartial^0_e + \alpha \in \mathscr{D}(E)\). Observe that given \(\xi \in \Gamma(M;E)\), \(\alpha(\xi)\in \Omega_\C^{(0,1)}(M;E)\), so the sum of \(\barpartial_e^0\) and \(\alpha\) makes sense as an operator. Moreover,
    \begin{align*}
        \parent{\barpartial_e^0 + \alpha}(f\cdot  \xi) & = \barpartial_e^0(f \cdot \xi) + \alpha(f \cdot \xi) = \barpartial f \otimes \xi + f \cdot \barpartial_e^0 \xi + f \cdot \alpha(\xi) \\
        & = \barpartial f \otimes \xi + f \cdot \parent{\barpartial_e^0+\alpha}(\xi).
    \end{align*}
    Reciprocally, choose another \(\barpartial_e \in \mathscr{D}(E)\). A direct computation shows that \(\barpartial_e^0-\barpartial_e\) is \(\mathscr{C}^\infty(M;\C)\)-linear, so \(\barpartial_e^0-\barpartial_e\in \Omega_\C^{(0,1)}(M;\mathrm{End}(E))\).
\end{proof}

Let \(\bar{\mathscr{A}}(E)\) be the set of all connections on a complex vector bundle over a complex manifold \(E \to M\) (again, we will simply write \(\bar{\mathscr{A}}\) if no confusion arises). Then, we have a map (see \eqref{nabla01_is_Dolbeault_operator}):
\[
\begin{array}{ccc}
    \bar{\mathscr{A}}(E) & \longrightarrow & \mathscr{D}(E) \\
    \nabla & \mapsto & \nabla^{(0,1)}.
\end{array}
\]
Suppose now that \(E\) is endowed with a hermitian fiber metric \(h\). This condition will allow us to give an inverse to the map we just defined. Let \(\mathscr{A}(E)\) be the set of all the connections on \(E\) that are compatible with the fiber metric. By applying \(\Pi^{(0,1)}\) to equation \ref{def_cond_compatibility}, we get
\[
\barpartial h(\xi, \eta) = h\parent{\nabla^{(0,1)}\xi, \eta} + h \parent{\xi, \nabla^{(1,0)}\eta}.
\]
This is key to define the inverse map. Indeed, we can define \(\mathscr{D}(E) \to \mathscr{A}(E)\) as follows: given a Dolbeault operator \( \barpartial_e\), we define a connection by \(\nabla:=\nabla^{(1,0)}_e + \barpartial_e\), where \(\nabla^{(0,1)}_e\colon \Gamma(M;E) \to \Omega_\C^{(1,0)}(M;E)\) is uniquely determined by
\begin{equation} \label{def_nabla_e}
\barpartial h(\xi, \eta) = h \parent{\barpartial_e \xi, \eta} + h\parent{\xi, \nabla^{(1,0)}_e \eta}.
\end{equation}
If the operator \(\nabla_e^{(0,1)}\) exists, it is unique, because any other operator \(\overline{\nabla}_e^{(1,0)}\) satisfying equation \eqref{def_nabla_e} should satisfy \(h\parent{\xi, \nabla_e^{(1,0)}\eta - \overline{\nabla}_e^{(1,0)}\eta}=0\) for every \(\xi, \eta \in \Gamma(M;E)\), thus \(\nabla_e^{(1,0)} =\overline{\nabla}_e^{(1,0)}\). Moreover, equation \eqref{def_nabla_e} implies that 
\begin{equation} \label{eq_ext_nabla01}
\nabla_e^{(1,0)} (f \cdot \xi) = \partial f \otimes \xi  + f \cdot \nabla_e^{(1,0)} \xi.
\end{equation}
Indeed, for any \(\eta, \xi \in \Gamma(M;E)\) one has
\begin{align*}
    h\parent{\xi, \nabla_e^{(1,0)} (f\cdot \eta)} & = \barpartial h(\xi, f\cdot \eta) - h\parent{\nabla^{(0,1)}\xi, f\cdot\eta} = \barpartial \parent{\bar{f}\cdot h(\xi, \eta)} + \bar{f}\cdot h\parent{\nabla^{(0,1)}\xi,\eta} \\
    & = \barpartial \bar{f} \cdot h(\xi, \eta) + \bar{f}\cdot \parent{h(\xi, \eta) + h\parent{\nabla^{(0,1)}\xi, \eta}} = \barpartial \bar{f} \cdot h(\xi, \eta) + \bar{f}\cdot h\parent{\xi, \nabla_e^{(1,0)}\eta} \\
    & = h\parent{\xi, \partial f \otimes \eta + f\cdot\nabla_e^{(1,0)}\eta}.
\end{align*}
This implies that \(\nabla\) is a well-defined connection. To check that the so-defined connection \(\nabla\) is compatible with the metric, observe that equation \eqref{def_nabla_e} implies
\begin{align*}
    \partial h(\xi, \eta) & = \overline{\barpartial \ \overline{h(\xi, \eta)}} = \overline{\barpartial h(\eta, \xi)} = \overline{h\parent{\barpartial_e \eta, \xi} + h\parent{\eta, \nabla_e^{(1,0)}\xi}} = h\parent{\xi, \barpartial_e \eta} + h\parent{\nabla_e^{(1,0)}\xi, \eta}.    
\end{align*}
Therefore, 
\[
\diff h(\xi, \eta) = \partial h(\xi, \eta) +  \barpartial h(\xi, \eta) = h\parent{\barpartial_e\xi + \nabla_e^{(1,0)}\xi, \eta} + h\parent{\xi, \barpartial_e \eta + \nabla^{(1,0)}_e \eta}.
\]
 To see that \(\nabla_e^{(1,0)}\) actually exists, we follow an argument similar to the one given in the proof of theorem \ref{th_Chern_connection}: we choose a local frame \((s_1, \ldots s_r)\) of \(E\) and define the operator on the frame. There exist \((0,1)\)-forms \(\alpha^r_i\) such that \(\barpartial_e s_i=\alpha^r_i \otimes s_r\) (this is analogous to remark \ref{coordinate_exp_connection}). Then, we can define \(\nabla_e^{(1,0)}\) by its action on a local frame \(\nabla_e^{(1,0)} (\xi^i \cdot s_i)=\parent{\partial \xi^i + \xi^j \cdot \omega ^k_j} \otimes s_k  \), where
 \[
 \omega^k_j := \overline{h^{mi}\cdot \parent{\barpartial h_{ij}-\alpha^r_i \cdot h_{rj}}}.
 \]
 Finally, observe that \(\nabla_e^{(0,1)}=\barpartial_e\) by definition. We thus have a bijection \(\mathscr{A}(E) \leftrightarrow \mathscr{D}(E)\).\\

 It is appropriate now to state (without proof) the following result.

 \Theorem{}{Let \(E \to M\) be a smooth complex vector bundle defined over a complex manifold \(X\). Let \(\barpartial_e\) be a Dolbeault operator satisfying the \textit{integrability condition} \(\barpartial_e \circ \barpartial_e =0\). Then, there is a unique holomorphic structure on \(E\) such that \(\barpartial_e=\barpartial_E\). \label{th_holomorphic_struc_iff_integrable}}
 \begin{proof}
    See \parencite[I-Prop. 3.7]{kobayashi_complex_vector_bundles}. Observe that by \(\barpartial_e^2=0\) we mean that the composition of \(\barpartial_e\) with its extension to one-forms \eqref{eq_ext_Dolb_operator} vanishes.
 \end{proof}
 Motivated by this result, let \(\mathscr{H}(E):=\corch{\barpartial_e \in \mathscr{D}(E): \barpartial_e^2=0}\). Observe that each element on \(\mathscr{H}(E)\) defines a holomorphic structure on \(E\). We now ask ourselves which connections correspond to \(\mathscr{H}(E)\) under the Chern correspondence. It turns out that this condition depends on the curvature. By \eqref{eq_decomp_E_forms}, we have a decomposition \(\Omega_\C^2(M;\mathrm{End}(E))=\Omega_\C^{(2,0)}(M; \mathrm{End}(E))\oplus \Omega_\C^{(1,1)}(M; \mathrm{End}(E)) \oplus \Omega_\C^{(0,2)}(M; \mathrm{End}(E))\), hence the curvature of any connection decomposes as \(F_\nabla=F_{\nabla}^{(2,0)}+F_\nabla^{(1,1)}+F_\nabla^{(0,2)}\).

 \Lemma{}{Let \(E \to M\) be a complex vector bundle, and let \(\nabla\) be a connection on \(E\). Then, \(\barpartial_e:=\nabla^{(0,1)}\) is an integrable Dolbeault operator if and only if \(F_\nabla^{(2,0)}=0\). \label{lemma_integrable_Dolb_operator}} 
 \begin{proof}
    First, observe that the decomposition of the curvature induces a decomposition of \(\diff_\nabla^2=\diff_\nabla^{(2,0)}+\diff_\nabla^{(1,1)}+\diff_\nabla^{(0,2)}\), where \(\diff_\nabla^{(0,2)}(\alpha)=F_\nabla^{(0,2)}\wedge \alpha\) for any \(\alpha \in \Omega_\C^{(p,q)}(M;E)\). The other operators are defined analogously. On the other hand, \(\nabla=\nabla^{(1,0)}+\barpartial_e\) means that, after extending these operators to act on higher degree forms as in \eqref{eq_ext_barpartial_e} and \eqref{eq_ext_nabla01}, we have
    \begin{align*}
    \diff_\nabla^2&=(\nabla^{(1,0)}+\barpartial_e) \circ (\nabla^{(1,0)}+\barpartial_e)\\
    &=\nabla^{(1,0)}\circ \nabla^{(1,0)} + \parent{\nabla^{(1,0)}\circ \barpartial_e+\barpartial_e\circ \nabla^{(1,0)}} + \barpartial_e \circ \barpartial_e. 
    \end{align*}
    Applying the projection \(\Pi^{(0,2)}\) to both expressions for \(\diff_\nabla^2\), we get \(\diff_\nabla^{(0,2)}=\barpartial_e^2\), and the result follows.
 \end{proof}
Observe that \(F_\nabla^{(0,2)}=0\) if and only if \(F_\nabla^{(2,0)}\). This follows from the fact that \(F_\nabla^{(2,0)}(\bar{X}, \bar{Y})=F_\nabla^{(0,2)}(X,Y)\) for all vector fields \(X, Y \in \Gamma(M; T_\C M)\). \\

Motivated by the previous lemma, we set \(\mathscr{A}^{(1,1)}(E):=\corch{\nabla \in \mathscr{A}: F_\nabla \in \Omega_\C^{(1,1)}(M; \mathrm{End}(E))}\). We call its elements connections of type \((1,1)\) or \textit{integrable connections}\index{Connection!integrable}. By the previous lemma, the Chern correspondence restricts to a bijection \(\mathscr{A}^{(1,1)}(E) \leftrightarrow \mathscr{H}(E)\). Observe that under this bijection, a connection becomes the Chern connection of the complex structure that its \((0,1)\) component induces on \(E\). \\

Looking at theorem \ref{th_holomorphic_struc_iff_integrable}, we would like to identify \(\mathscr{H}(E)\) with all possible holomorphic structures on \(E\) (whose underlying smooth complex vector bundle structure is that of \(E\)). Although theorem \ref{th_holomorphic_struc_iff_integrable} does state that every element \(\barpartial_e \in \mathscr{H}(E)\) induces a holomorphic structure on \(E\), it may happen that two different integrable Dolbeault operators induce biholomorphic complex structures, and it would not make sense to count those as different. To remedy this issue, let us define the following groups.

\Definition{}{Let \(E \to M\) be a complex vector bundle endowed with a hermitian fiber product \(h\). The \textit{complex gauge group} \(\mathscr{G}^\C=\mathrm{GL}(E) \subset \Gamma(M; \mathrm{End}(E))\) is the group of complex smooth bundle automorphisms (equipped with fiberwise composition). The \textit{gauge group} or \textit{group of unitary automorphisms} \(\mathscr{G}=\mathrm{U}(E,h)\subset \mathscr{G}^\C\) is the group of bundle isomorphisms of \(E\) satisfying \(h_p(f_p(\xi), f_p(\eta))=h_p(\xi, \eta)\), \(\xi, \eta \in E_p\) for all \(p \in M\).}

The complex gauge group acts through the right on the space of Dolbeault operators as follows 
\[
\begin{array}{ccc}
    \mathscr{D}(E) \times \mathscr{G}^\C & \longrightarrow & \mathscr{D}(E) \\
    (\barpartial_e, f) & \mapsto & \barpartial_e \cdot f
\end{array}
\]
where \(\barpartial_e \cdot f:=\parent{\mathrm{id}_{T_\C^*M}\otimes f^{-1}} \circ \barpartial_e \circ f\). The identity component is usually omitted, and one writes \(f^{-1} \circ \barpartial_e \circ f\). 

\Proposition{}{Let \(E \to M\) be a complex vector bundle over a smooth manifold. Let \(f \in \mathscr{G}^\C\), and \(\barpartial_e, \barpartial_e'\) two Dolbeault operators such that \(\barpartial_e'=\barpartial_e\cdot f\). Let \(\mathcal{E}\) and \(\mathcal{E}'\) be the complex structures induced on \(E\) by theorem \ref{th_holomorphic_struc_iff_integrable}. Then, \(f\colon \mathcal{E' \to \mathcal{E}}\) is an isomorphism between holomorphic vector bundles.}
\begin{proof}
    Since \(f \in \mathscr{G}^\C\), it is an isomorphism of complex smooth vector bundles, so one needs only check that it is holomorphic. This follows from the fact that given two holomorphic vector bundles, holomorphicity of bundle homomorphism is determined by the condition \(f \circ \barpartial_e'=\barpartial_e\circ f\).
\end{proof}

Thus, the true set of complex structures is given by the quotient of \(\mathscr{H}(E)\) by the action of the complex gauge group:
\[
\mathscr{C}(E)=\frac{\mathscr{H}(E)}{\mathscr{G}^\C},
\]
\(\mathscr{C}\) being the set of complex structures on \(E\). Moreover, \(\mathscr{G}^\C\) also acts on \(\mathscr{A}(E)\) by \(\nabla \cdot f:=f^* \circ \nabla^{(1,0)} \circ (f^*)^{-1} + f^{-1} \circ \nabla^{(0,1)} \circ f\) (same abuse of notation as before), where \(f^*\) is the adjoint of \(f\), i.e. the unique bundle homomorphism satisfying \(h_p(\xi, f_p^*(\eta))=h_p(f_p(\xi), \eta)\) for \(\xi, \eta \in E_p\) and all \(p \in M\). One may show (see \parencite[Sec. 7.1]{kobayashi_complex_vector_bundles}) that this action maps \(\mathscr{A}^{(1,1)}(E)\) to itself, so it restricts to an action on \(\mathscr{A}^{(1,1)}(E)\). We then get a bijection 
\begin{equation} \label{eq_moduli_Chern_correspondence}
\frac{\mathscr{A}^{(1,1)}(E)}{\mathscr{G}^\C} \leftrightarrow \frac{\mathscr{H}(E)}{\mathscr{G}^\C} = \mathscr{C}(E).
\end{equation}
These spaces admit extra structures (i.e. topologies and even manifold structures). When endowed with those, we call them \textit{Moduli spaces}.